\documentclass[border=5pt]{standalone}
\usepackage{mathptmx}
\usepackage[T1]{fontenc}
\usepackage{graphicx}
\usepackage{tikz}
\usetikzlibrary{shapes.geometric,shapes.misc,arrows.meta,positioning,calc,fit,shadows.blur}
\usepackage{xcolor}

\definecolor{diagLaw}{HTML}{1D6FE0}
\definecolor{diagPlan}{HTML}{E63946}
\definecolor{diagEnv}{HTML}{12A150}
\definecolor{diagLatent}{HTML}{7C3AED}
\definecolor{diagProbe}{HTML}{E8890C}
\definecolor{gapblue}{HTML}{00B4D8}
\definecolor{gapred}{HTML}{C2185B}
\definecolor{inkmute}{HTML}{6B7280}

\begin{document}
\begin{tikzpicture}[
    >=Stealth,
    font=\sffamily\small,
    % Flat fill, heavy tinted border, whisper of a shadow. No gradients: at print size they eat
    % text contrast. Rounded rectangles align on their edges, which ellipses cannot -- that is what
    % lets the three bands below actually line up.
    card/.style={draw, line width=1.6pt, rounded corners=3.5pt, align=center, fill=white,
                 minimum height=9mm, inner xsep=2.6mm, font=\sffamily\small\bfseries,
                 blur shadow={shadow blur steps=8, shadow xshift=0.5pt, shadow yshift=-0.8pt,
                              shadow opacity=11, shadow blur radius=1.3pt}},
    pic/.style={card, inner sep=1.1pt, minimum height=0pt},   % image cards: no text padding
    env/.style ={card, draw=diagEnv!70!black,    text=diagEnv!25!black},
    lat/.style ={card, draw=diagLatent!70!black, text=diagLatent!25!black},
    prb/.style ={card, draw=diagProbe!75!black,  text=diagProbe!30!black},
    law/.style ={card, draw=diagLaw!70!black,    text=diagLaw!25!black},
    pln/.style ={card, draw=diagPlan!70!black,   text=diagPlan!25!black},
    flow/.style={->, line width=0.9pt, draw=inkmute!65},
    % THE MOTIF: the two gaps are the conceptual punchline, so they are the heaviest marks here.
    gap/.style={line width=2.6pt, dash pattern=on 6pt off 3pt, -{Stealth[length=7pt]}},
    caps/.style={fill=#1, text=white, rounded corners=4pt, inner xsep=5pt, inner ysep=2.4pt,
                 font=\sffamily\scriptsize\bfseries},
    bandlab/.style={font=\sffamily\scriptsize\bfseries, text=inkmute!85, anchor=east, align=right},
    plab/.style={font=\sffamily\scriptsize, text=inkmute, align=center},
]

%=========== ROW 1 : the legal source, and the CLASSICAL isomorphism ===========%
\node[law] (statute) at (2.6,3.15)  {Legal source text};
\node[law] (rules)   at (8.4,3.15)  {DDL rule base};
\draw[gap, draw=inkmute!65, line width=1.9pt, dash pattern=on 3.5pt off 2.5pt]
      (statute) -- node[caps=inkmute!70] {isomorphism} (rules);

%=========== ROW 2 : world -> atoms  (the GROUNDING gap) ===========%
\node[pic, draw=diagEnv!70!black] (world)
      {\includegraphics[width=15mm]{inset_world.png}};
\node[plab, below=0.8mm of world, text=diagEnv!30!black] {\bfseries World};
\node[lat] (wm)     at (4.0,0.55) {World model $\psi$};
\node[prb] (probes) at (7.7,0.55) {Probes $\phi$};
\node[law] (atoms)  at (12.4,0.55) {Atoms};
\draw[flow] (world) -- (wm);
\draw[flow] (wm) -- (probes);
\draw[gap, draw=gapblue] (probes) -- node[caps=gapblue] {Grounding gap} (atoms);

%=========== ROW 3 : verdict -> constraint  (the ONTOLOGICAL gap) ===========%
\node[law] (verdict) at (12.4,-2.5) {Verdict};
\node[pln] (cons)    at (7.7,-2.5)  {Planning constraint};
\node[pic, draw=diagPlan!70!black] (planner) at (3.1,-2.5)
      {\includegraphics[width=15mm]{inset_planner.png}};
\node[plab, below=0.8mm of planner, text=diagPlan!25!black] {\bfseries RRT planner};

\draw[flow] (rules.east) -- ++(1.1,0) |- (atoms.east);
\draw[flow] (atoms) -- (verdict);
\draw[gap, draw=gapred] (verdict) -- node[caps=gapred] {Ontological gap} (cons);
\draw[flow] (cons) -- (planner);

%=========== the closed loop, made explicit ===========%
\draw[flow, line width=1.1pt] (planner.west) -- ++(-1.5,0)
      node[midway, above=1pt, font=\sffamily\scriptsize, text=inkmute] {action} |- (world.west);

%=========== band labels ===========%
\node[bandlab] at (0.55,3.15)  {\textsc{law}};
\node[bandlab] at (0.55,1.5)   {\textsc{perception}};
\node[bandlab] at (0.55,-1.35) {\textsc{control}};

\end{tikzpicture}
\end{document}
